\documentclass[11pt]{article}

\usepackage[T1]{fontenc}
\usepackage[utf8]{inputenc}
\usepackage{lmodern}
\usepackage[a4paper,margin=27mm]{geometry}
\usepackage{amsmath,amssymb,amsthm,mathtools}
\usepackage{microtype}
\usepackage{xcolor}
\usepackage{enumitem}
\usepackage{booktabs}
\usepackage{tabularx}
\usepackage{hyperref}

\hypersetup{
  colorlinks=true,
  linkcolor=blue!55!black,
  citecolor=blue!55!black,
  urlcolor=blue!55!black
}

\newtheorem{theorem}{Theorem}[section]
\newtheorem{proposition}[theorem]{Proposition}
\newtheorem{lemma}[theorem]{Lemma}
\newtheorem{corollary}[theorem]{Corollary}
\newtheorem{definition}[theorem]{Definition}
\theoremstyle{remark}
\newtheorem{remark}[theorem]{Remark}

\newcommand{\A}{\mathcal A}
\newcommand{\B}{\mathcal B}
\newcommand{\Hh}{\mathcal H}
\newcommand{\C}{\mathbb C}
\newcommand{\R}{\mathbb R}
\newcommand{\Z}{\mathbb Z}
\newcommand{\id}{\mathbf 1}
\newcommand{\Tr}{\operatorname{Tr}}
\newcommand{\Ad}{\operatorname{Ad}}
\newcommand{\Ind}{\operatorname{Ind}}
\newcommand{\CAR}{\operatorname{CAR}}
\newcommand{\Hom}{\operatorname{Hom}}
\newcommand{\KMS}{\mathrm{KMS}}
\newcommand{\Fix}{\operatorname{Fix}}
\newcommand{\dd}{\,\mathrm d}

\title{The Holomorphic Rotational KMS Extension Theorem:\\
Index Four, the Canonical Expectation, and the TFPT Finite Shadow}
\author{Proof memorandum}
\date{29 August 2026}

\begin{document}
\maketitle

\begin{abstract}
Let $\A\subset\B$ be an irreducible finite-index inclusion of conformal
nets on $S^1$, let $\A=\B^{\Z_4}$ for a proper gauge action commuting
with rotations, and suppose that $\B$ is holomorphic, with
$\mu_{\B}=1$ and central charge $8$.  We prove the strongest correct
form of the proposed KMS extension statement.  At angular inverse
temperature $\beta>0$, the vacuum-sector Gibbs state of $\B$ restricts
to a rotational KMS state $\omega_\beta$ of $\A$.  The group average
\[
 E=\frac14\sum_{r=0}^{3}\alpha_r
\]
is the unique normal faithful conditional expectation
$\B(I)\to\A(I)$ on every interval, and the extension is
$\widetilde\omega_\beta=\omega_\beta\circ E$.  It is the unique
$\Z_4$-invariant state extension, even before the KMS condition is
imposed.  Finally, the Kawahigashi--Longo--M\"uger sector theorem
together with the rotational KMS classification of
Longo--Tanimoto implies that a holomorphic conformal net has only one
rotational $\beta$-KMS state.  Hence every rotational KMS extension is
automatically $\Z_4$-invariant and equals
$\omega_\beta\circ E$.

Two qualifications are essential.  The vacuum \emph{vector} state is
a rotation ground state, not a finite-temperature rotational KMS
state; ``vacuum KMS'' below means the vacuum-\emph{sector Gibbs
state}.  Also, $\mu=1$ eliminates DHR sectors but does not by itself
classify thermal states: the Longo--Tanimoto theorem is the necessary
analytic bridge.  We quote in full the precise forms of Takesaki's
expectation theorem, the finite-index/Q-system theorem, the KLM
$\mu$-index theorem, Araki's CAR KMS theorem, and the rotational KMS
classification, and verify their hypotheses against the frozen TFPT
finite-collar data.  The final section isolates the remaining
continuum identification in one box.
\end{abstract}

\tableofcontents

\section{Setting and correction of terminology}

\subsection{Conformal-net hypotheses}

Let $I\mapsto\B(I)$ be an irreducible local conformal net on $S^1$ in
its vacuum representation $(\Hh_{\B},U_{\B},\Omega_{\B})$.  Let
$G=\Z_4$ act properly by vacuum-preserving internal symmetries
$\alpha_r$, meaning
\[
 \alpha_r(\B(I))=\B(I),\qquad
 \alpha_r\Ad U_{\B}(g)=\Ad U_{\B}(g)\alpha_r
\]
for every interval $I$, every conformal transformation $g$, and
$r\in\Z_4$.  Put
\[
 \A(I)=\B(I)^G.
\]
We assume that the local inclusion is irreducible and has Jones index
four:
\[
 [\B(I):\A(I)]=4,\qquad
 \A(I)'\cap\B(I)=\C\id.
\tag{1.1}\label{eq:local-inclusion}
\]
For a proper outer action of a finite group on a factor, the first
identity follows from the fixed-point theorem; it is also included
explicitly among our hypotheses so no outerness convention is hidden.

The net $\B$ is assumed holomorphic in the operator-algebraic sense:
it is completely rational and
\[
 \mu_{\B}=1.
\tag{1.2}\label{eq:holomorphic}
\]
The central charge $c=8$ identifies the intended endpoint as the
$(E_8)_1$ net, but is not needed for KMS uniqueness once
\eqref{eq:holomorphic} is known.

\subsection{The rotational KMS state}

Let $L_0^{\B}\geq0$ be the conformal Hamiltonian and let
$\tau_t=\Ad e^{itL_0^{\B}}$ be the rotation dynamics on the locally
normal universal algebra $C^*_{\mathrm{ln}}(\B)$.  For $\beta>0$ put
\[
 \Phi_\beta^{\B}(x)
 =\frac{\Tr_{\Hh_{\B}}(e^{-\beta L_0^{\B}}x)}
 {\Tr_{\Hh_{\B}}(e^{-\beta L_0^{\B}})}.
\tag{1.3}\label{eq:gibbs}
\]
For a completely rational conformal net the character is finite, and
the rotational KMS classification quoted below in fact recovers the
trace-class property from any extremal KMS state.  In the concrete
$(E_8)_1$ case it is also immediate from the character
$E_4/\eta^8$.

\begin{definition}[Geometric rotational state used here]
\label{def:geometric}
The state on the subnet is
\[
 \omega_\beta:=\Phi_\beta^{\B}|_{C^*_{\mathrm{ln}}(\A)}.
\tag{1.4}\label{eq:omega-def}
\]
At the canonical angular normalization we set
$\beta=\beta_{\mathrm{angle}}=2\pi$.
\end{definition}

\begin{remark}[Ground state versus rotational KMS state]
The vector functional
$\langle\Omega_{\B},\,\cdot\,\Omega_{\B}\rangle$ is a ground state for
rotations.  It is not the finite-temperature state
\eqref{eq:gibbs}.  The geometric KMS state on the real line for
translations/dilations and the rotational Gibbs state on the circle
are likewise different constructions.  The theorem below concerns
the latter, because its unconditional classification by
superselection sectors is available in exactly this setting.
\end{remark}

\section{Classical theorems used, stated in full}

\subsection{Araki's CAR KMS theorem}

\begin{theorem}[Araki: unique quasi-free CAR equilibrium state]
\label{thm:araki}
Let $\mathfrak h$ be a complex Hilbert space, let $h=h^*$ be a
possibly unbounded one-particle Hamiltonian, and let
$\tau_t^h$ be the strongly continuous Bogoliubov dynamics on
$\CAR(\mathfrak h)$ determined on generators by
\[
 \tau_t^h(a(f))=a(e^{ith}f).
\]
For every $\beta>0$ there is a gauge-invariant quasi-free state
$\varphi_{\beta,h}$ whose two-point function is
\[
 \varphi_{\beta,h}(a^*(f)a(g))
 =\langle g,C_{\beta,h}f\rangle,\qquad
 C_{\beta,h}=(\id+e^{\beta h})^{-1},
\tag{2.1}\label{eq:araki-cov}
\]
and whose higher even correlation functions are the determinants
prescribed by the CAR Wick rule, while odd correlation functions
vanish.  This state is a $\beta$-KMS state for $\tau^h$.  In the
factorial self-dual CAR representation (in particular when the
covariance has no $0$ or $1$ eigenvalue in the finite-dimensional
case), it is the unique $\beta$-KMS state for that Bogoliubov
dynamics.  Thus the KMS state is completely determined by its
two-point covariance.
\end{theorem}

This is the KMS uniqueness theorem in
H.~Araki, \emph{On quasifree states of CAR and Bogoliubov
automorphisms}, Publ.\ Res.\ Inst.\ Math.\ Sci.\ \textbf{6}
(1970/71), 385--442, Theorem~3 and Eq.~(6.9), translated from the
self-dual notation to the ordinary gauge-invariant CAR notation.
Araki's factor proviso is recorded rather than suppressed.

\subsection{Takesaki's conditional-expectation theorem}

\begin{theorem}[Takesaki: modular invariance criterion]
\label{thm:takesaki}
Let $M$ be a von Neumann algebra, let $\phi$ be a faithful normal
semifinite weight on $M$, and let $N\subset M$ be a von Neumann
subalgebra such that $\phi|_N$ is semifinite.  Write
$\sigma^\phi=(\sigma_t^\phi)_{t\in\R}$ for the modular automorphism
group of $\phi$.  The following are equivalent:
\begin{enumerate}[label=\textup{(\roman*)}]
\item $N$ is globally invariant under the modular group,
  $\sigma_t^\phi(N)=N$ for every $t\in\R$;
\item there exists a normal faithful conditional expectation
  $E_\phi:M\to N$ such that
  $\phi\circ E_\phi=\phi$ on $M_+$.
\end{enumerate}
When these conditions hold, $E_\phi$ is unique among the normal
conditional expectations preserving $\phi$, it commutes with the
modular group,
\[
 E_\phi\sigma_t^\phi=\sigma_t^{\phi|_N}E_\phi,
\tag{2.2}\label{eq:takesaki-intertwine}
\]
and the modular group of the restricted weight is the restriction:
\[
 \sigma_t^{\phi|_N}=\sigma_t^\phi|_N.
\tag{2.3}\label{eq:mod-restrict}
\]
\end{theorem}

This is M.~Takesaki, \emph{Conditional expectations in von Neumann
algebras}, J.\ Funct.\ Anal.\ \textbf{9} (1972), 306--321, and
\emph{Theory of Operator Algebras II}, Springer, 2003,
Theorem~IX.4.2.

\subsection{Finite index, Radon--Nikodym uniqueness, and Q-systems}

\begin{theorem}[Finite-index expectation and Q-system package]
\label{thm:qsystem}
Let $N\subset M$ be an inclusion of infinite factors with finite
index and separable predual.
\begin{enumerate}[label=\textup{(\alph*)}]
\item There is a normal faithful conditional expectation
  $E_0:M\to N$ of finite index.  Fix one such expectation.  Every
  other normal faithful conditional expectation $F:M\to N$ is of
  the form
  \[
    F(x)=E_0(h^{1/2}xh^{1/2}),\qquad x\in M,
  \tag{2.4}\label{eq:rn-expectation}
  \]
  for a positive invertible $h\in N'\cap M$ satisfying
  $E_0(h)=\id$.  Conversely every such $h$ defines an expectation.
  In particular, if $N'\cap M=\C\id$, then the normal faithful
  conditional expectation is unique.
\item Let $\iota:N\hookrightarrow M$ be the inclusion morphism and
  $\bar\iota:M\to N$ a conjugate.  The dual canonical endomorphism
  $\theta=\bar\iota\iota\in\operatorname{End}(N)$, together with
  standard conjugacy intertwiners
  \[
    w\in\Hom(\id_N,\theta),\qquad
    x\in\Hom(\theta,\theta^2),
  \]
  is a standard Q-system: $w$ is a unit, $x$ is associative and
  unital with $w$, and $x$ obeys the Frobenius and standardness
  relations.  Conversely, every standard Q-system in
  $\operatorname{End}(N)$ reconstructs, up to isomorphism fixing
  $N$, a finite-index extension $N\subset M$.  Its dimension is
  \[
    d(\theta)=[M:N].
  \tag{2.5}\label{eq:q-index}
  \]
  If the endomorphisms lie in the braided DHR category of a
  conformal net, the reconstructed extension is local exactly when
  the Q-system is commutative with respect to the braiding.
\item If $M$ carries an outer action of a finite group $G$ and
  $N=M^G$, then
  \[
    E_0(x)=\frac1{|G|}\sum_{g\in G}\alpha_g(x),\qquad
    [M:N]=|G|.
  \tag{2.6}\label{eq:group-average}
  \]
  The dual canonical endomorphism is the regular Q-system.  For
  $G=\Z_4$ it is the sum of the four one-dimensional character
  sectors and has dimension four.
\end{enumerate}
\end{theorem}

Part (a) is the Radon--Nikodym description of expectations in the
finite-index theory of H.~Kosaki and R.~Longo; see R.~Longo,
\emph{Index of subfactors and statistics of quantum fields I},
Commun.\ Math.\ Phys.\ \textbf{126} (1989), 217--247, especially
the uniqueness/minimal-index discussion, and
F.~Hiai, \emph{Minimum index for subfactors and entropy II},
J.\ Math.\ Soc.\ Japan \textbf{43} (1991), 347--379.  Parts (b) and
(c) are the subfactor/Q-system and fixed-point forms used in
R.~Longo and K.-H.~Rehren, \emph{Nets of subfactors},
Rev.\ Math.\ Phys.\ \textbf{7} (1995), 567--597, together with the
local-extension formulation of J.~B\"ockenhauer and D.~E.~Evans,
\emph{Modular invariants, graphs and $\alpha$-induction for nets of
subfactors I}, Commun.\ Math.\ Phys.\ \textbf{197} (1998), 361--386.

\subsection{\texorpdfstring{The KLM $\mu$-index theorem}{The KLM mu-index theorem}}

\begin{theorem}[Kawahigashi--Longo--M\"uger]
\label{thm:klm}
Let $\mathcal N$ be a local irreducible conformal net satisfying the
split property and strong additivity, and suppose that its
two-interval inclusion
\[
 \mathcal N(E)\subset\widehat{\mathcal N}(E)
 :=\mathcal N(E')'
\]
has finite index.  Then $\mathcal N$ is completely rational; it has
finitely many inequivalent irreducible DHR sectors
$\rho_0=\id,\rho_1,\ldots,\rho_n$, all of finite statistical
dimension, every locally normal representation is a direct sum of
these sectors, the braiding is non-degenerate, and
\[
 \mu_{\mathcal N}
 =[\widehat{\mathcal N}(E):\mathcal N(E)]
 =\sum_{j=0}^{n}d(\rho_j)^2.
\tag{2.7}\label{eq:mu-sum}
\]
The two-interval inclusion is the Longo--Rehren inclusion associated
with the full sector category.  If
$\mathcal N\subset\mathcal M$ is a finite-index irreducible local
extension of completely rational nets, then
\[
 \mu_{\mathcal M}
 =\frac{\mu_{\mathcal N}}{[\mathcal M:\mathcal N]^2}.
\tag{2.8}\label{eq:mu-extension}
\]
Consequently, $\mu_{\mathcal M}=1$ if and only if the vacuum is the
only irreducible DHR sector of $\mathcal M$.
\end{theorem}

This is Y.~Kawahigashi, R.~Longo and M.~M\"uger,
\emph{Multi-interval subfactors and modularity of representations in
conformal field theory}, Commun.\ Math.\ Phys.\ \textbf{219} (2001),
631--669: the two-interval/Longo--Rehren identification, the equality
of $\mu$ with the global sector dimension, and the extension-index
formula are the results used here.  The final equivalence follows
from \eqref{eq:mu-sum}, since $d(\rho_0)=1$ and every nonzero
statistical dimension is at least one.

\subsection{Rotational KMS classification}

\begin{theorem}[Longo--Tanimoto]
\label{thm:lt}
Let $\mathcal N$ be a completely rational conformal net and let
$C^*_{\mathrm{ln}}(\mathcal N)$ be its locally normal universal
$C^*$-algebra.  Then
\[
 C^*_{\mathrm{ln}}(\mathcal N)
 \cong\bigoplus_{j=0}^{n}B(\Hh_j),
\tag{2.9}\label{eq:universal-sum}
\]
where the summands correspond bijectively to the inequivalent
irreducible DHR representations $\rho_j$.  Every extremal rotational
$\beta$-KMS state is the Gibbs state in one irreducible
representation:
\[
 \Phi_{j,\beta}(x)
 =\frac{\Tr_{\Hh_j}(e^{-\beta L_0^{(j)}}\rho_j(x))}
 {\Tr_{\Hh_j}(e^{-\beta L_0^{(j)}})}.
\tag{2.10}\label{eq:sector-gibbs}
\]
In particular $e^{-\beta L_0^{(j)}}$ is trace class for every sector
that supports such a state.  Every rotational $\beta$-KMS state is a
convex combination of the extremal sector Gibbs states.  Hence, if
the vacuum is the only irreducible sector, the vacuum-sector Gibbs
state is the unique rotational $\beta$-KMS state.
\end{theorem}

Equation \eqref{eq:universal-sum} is the completely rational universal
algebra theorem of Carpi--Conti--Hillier--Weiner used in
R.~Longo and Y.~Tanimoto, \emph{Rotational KMS states and type I
conformal nets}, Commun.\ Math.\ Phys.\ \textbf{357} (2018),
249--266, Theorems~3.1, 3.2 and 3.7.  The last sentence is an
immediate specialization to a one-summand algebra.

\section{Verification of the hypotheses}

\subsection{The index-four fixed-point Q-system}

For each interval set $N=\A(I)$ and $M=\B(I)$.  Conformal-net local
algebras are factors; \eqref{eq:local-inclusion} gives finite index
and trivial relative commutant.  Theorem~\ref{thm:qsystem} therefore
applies.  The group average is
\[
 E_I(x)=\frac14\bigl(x+\alpha_1(x)+\alpha_2(x)+\alpha_3(x)\bigr).
\tag{3.1}\label{eq:EI}
\]
It is normal, unital, positive, faithful, idempotent, and
$\A(I)$-bimodular.  For example, if $a,b\in\A(I)$, then
\[
 E_I(axb)=\frac14\sum_{r=0}^{3}a\alpha_r(x)b
 =aE_I(x)b.
\tag{3.2}\label{eq:bimodule}
\]
Faithfulness follows because $x\geq0$ and $E_I(x)=0$ imply that the
sum of four positive operators is zero, hence $x=0$.  Compatibility
under interval inclusion is immediate from the same formula, so the
$E_I$ define a net expectation and a global group average
\[
 E:C^*_{\mathrm{ln}}(\B)\longrightarrow
 C^*_{\mathrm{ln}}(\A).
\tag{3.3}\label{eq:global-E}
\]

The regular $\Z_4$ Q-system has
\[
 \theta\simeq\rho_0\oplus\rho_1\oplus\rho_2\oplus\rho_3,
\qquad d(\rho_r)=1,\qquad d(\theta)=4.
\tag{3.4}\label{eq:regular-q}
\]
The frozen TFPT simple-current arithmetic
$h(J^r)=1$ for $r=1,2,3$ verifies the integer-spin/locality condition
for the intended index-four extension.  The arithmetic cited in
\texttt{v973} is
\[
 \mu_{\A}=16,\qquad [\B:\A]=4,\qquad
 \mu_{\B}=\frac{16}{4^2}=1.
\tag{3.5}\label{eq:tfpt-mu}
\]
This is precisely the KLM extension formula
\eqref{eq:mu-extension}; the Python computation is a finite exact
check of its numerical inputs, not a replacement for the theorem.

\subsection{The frozen CAR covariance}

The finite collar used by \texttt{v995} has eight complex fermion
modes, hence
\[
 \CAR_8\simeq M_{2^8}(\C)=M_{256}(\C).
\]
In dimensionless modular units the one-particle Hamiltonian is
\[
 h_{\mathrm{collar}}
 =\Delta\,\operatorname{diag}(1,2,\ldots,8),
\qquad
 \Delta=6\log(3/2),
\tag{3.6}\label{eq:collar-h}
\]
and the script sets $\beta=1$.  Therefore
\[
 q=e^{-\Delta}=(2/3)^6=\frac{64}{729},
\qquad
 C_j=\frac{q^j}{1+q^j}
 =\frac1{1+e^{j\Delta}},\quad j=1,\ldots,8.
\tag{3.7}\label{eq:frozen-C}
\]
This is exactly Araki's covariance
\eqref{eq:araki-cov}.  All eigenvalues lie strictly between $0$ and
$1$, so the finite CAR representation is factorial and the density
matrix is faithful.  On the occupation basis
$n=(n_1,\ldots,n_8)\in\{0,1\}^8$,
\[
 p(n)=Z^{-1}q^{\sum_{j=1}^{8}jn_j},
\qquad
 Z=\prod_{j=1}^{8}(1+q^j).
\tag{3.8}\label{eq:gibbs-weights}
\]
Thus all $256$ weights are positive, and
\[
 \frac{p(n+e_j)}{p(n)}=q^j
\]
whenever $n_j=0$.  These are the exact detailed-balance relations
checked in \texttt{v995}.  The angular convention
$\beta_{\mathrm{angle}}=2\pi$ is recovered by absorbing $2\pi$ into
the modular Hamiltonian: $\beta h=(2\pi)H_{\mathrm{angle}}$.

\begin{remark}[What Araki proves here]
Araki's theorem proves uniqueness of the KMS state of the full CAR
dynamics and its determination by \eqref{eq:frozen-C}.  Restricting
that state to the even/bilinear $\mathrm{SO}(16)_1$ subnet selects the
vacuum-sector thermal state.  It does \emph{not} say that the
observable $\mathrm{SO}(16)_1$ net has only one rotational KMS state:
that net has four DHR sectors and hence four extremal sector Gibbs
states.  Holomorphic extension to $\B$ is what removes this sector
choice.
\end{remark}

\subsection{Holomorphy and thermal classification}

By \eqref{eq:holomorphic} and Theorem~\ref{thm:klm},
\[
 1=\mu_{\B}
 =\sum_{\rho\in\operatorname{Irr}(\B)}d(\rho)^2.
\tag{3.9}\label{eq:one-sector}
\]
The vacuum already contributes one, so there can be no other
irreducible DHR sector.  Theorem~\ref{thm:lt} then has one Gibbs
state and no other extreme point.  This verifies the analytic
hypothesis needed for unconditional rotational KMS uniqueness.

\section{The holomorphic KMS extension theorem}

\begin{theorem}[Holomorphic rotational KMS extension]
\label{thm:main}
Let $\A\subset\B$ satisfy the hypotheses of Section~1: for every
interval $I$, $\A(I)\subset\B(I)$ is an irreducible index-four
inclusion, $\A=\B^{\Z_4}$ for a proper internal gauge action
commuting with rotations, and $\B$ is a holomorphic conformal net
with $\mu_{\B}=1$ and $c=8$.  Let $\beta>0$ and let
$\omega_\beta$ be the restriction \eqref{eq:omega-def} of the
vacuum-sector rotational Gibbs state.  Then:
\begin{enumerate}[label=\textup{(\alph*)}]
\item $\omega_\beta$ has a rotational $\beta$-KMS extension to
  $\B$.
\item The maps
  \[
    E_I=\frac14\sum_{r=0}^{3}\alpha_r:
    \B(I)\to\A(I)
  \]
  are the unique normal faithful conditional expectations.  They
  form a net expectation $E$, and
  \[
    \widetilde\omega_\beta
    =\omega_\beta\circ E
    =\Phi_\beta^{\B}.
  \tag{4.1}\label{eq:main-extension}
  \]
  This is the unique $\Z_4$-invariant state extension of
  $\omega_\beta$, hence in particular the unique
  $\Z_4$-invariant rotational KMS extension.
\item Every rotational $\beta$-KMS extension of $\omega_\beta$ is
  $\Z_4$-invariant.  Consequently
  $\widetilde\omega_\beta$ is the unique rotational KMS extension
  without an invariance qualifier.
\end{enumerate}
\end{theorem}

\begin{proof}
\emph{Part (a).}
The Gibbs functional \eqref{eq:gibbs} is a rotational
$\beta$-KMS state on $C^*_{\mathrm{ln}}(\B)$ by
Theorem~\ref{thm:lt}, or directly by the cyclicity of the trace on
analytic elements:
\[
\begin{aligned}
 \Phi_\beta^{\B}(x\tau_{t+i\beta}(y))
 &=Z_\beta^{-1}
   \Tr(e^{-\beta L_0}x e^{itL_0}e^{-\beta L_0}
        y e^{\beta L_0}e^{-itL_0})\\
 &=Z_\beta^{-1}
   \Tr(e^{-\beta L_0}\tau_t(y)x)
 =\Phi_\beta^{\B}(\tau_t(y)x).
\end{aligned}
\tag{4.2}\label{eq:gibbs-kms}
\]
Its restriction to a $\tau$-invariant $C^*$-subalgebra remains KMS:
the same strip-analytic functions work for $x,y\in
C^*_{\mathrm{ln}}(\A)$.  By definition that restriction is
$\omega_\beta$.  Hence $\Phi_\beta^{\B}$ is the required extension.

\emph{Part (b): expectation and its uniqueness.}
Section~3.1 proves directly that the group average $E_I$ is a normal
faithful conditional expectation.  If $F_I:\B(I)\to\A(I)$ is any
other normal faithful conditional expectation, the
Radon--Nikodym part of Theorem~\ref{thm:qsystem} gives
\[
 F_I(x)=E_I(h^{1/2}xh^{1/2})
\]
with $h\in\A(I)'\cap\B(I)$.  Irreducibility gives
$h=\lambda\id$.  Since $F_I(\id)=\id=E_I(\id)$, $\lambda=1$, and
$F_I=E_I$.  This also shows why the bimodule property in the
definition of a conditional expectation is load-bearing: without
it, the Radon--Nikodym theorem does not apply, exactly as the
finite mutants in \texttt{v993} demonstrate.

The Gibbs state is gauge invariant.  Indeed the internal symmetry
commutes with $L_0$ and is unitarily implemented, so
\[
\begin{aligned}
 \Phi_\beta^{\B}(\alpha_r(x))
 &=Z_\beta^{-1}\Tr(e^{-\beta L_0}V_rxV_r^*)\\
 &=Z_\beta^{-1}\Tr(V_r^*e^{-\beta L_0}V_rx)
 =\Phi_\beta^{\B}(x).
\end{aligned}
\tag{4.3}\label{eq:gauge-invariance}
\]
Consequently
\[
 \Phi_\beta^{\B}(E(x))
 =\frac14\sum_{r=0}^{3}\Phi_\beta^{\B}(\alpha_r(x))
 =\Phi_\beta^{\B}(x).
\tag{4.4}\label{eq:E-preserves}
\]
Because $E(x)\in\A$ and
$\Phi_\beta^{\B}|_{\A}=\omega_\beta$,
\[
 \Phi_\beta^{\B}(x)
 =\Phi_\beta^{\B}(E(x))
 =\omega_\beta(E(x)),
\]
which proves \eqref{eq:main-extension}.

Now let $\psi$ be any $\Z_4$-invariant state on $\B$ whose restriction
to $\A$ is $\omega_\beta$.  No KMS hypothesis is needed:
\[
 \psi(x)
 =\frac14\sum_{r=0}^{3}\psi(\alpha_r(x))
 =\psi(E(x))
 =\omega_\beta(E(x)).
\tag{4.5}\label{eq:inv-unique}
\]
Thus $\psi=\omega_\beta\circ E$.

\emph{The modular-theory chain.}
We spell out the Takesaki step because it is frequently reversed
incorrectly.  In the GNS representation of the faithful KMS state,
the modular group and the physical dynamics are related, with the
present KMS convention, by
\[
 \sigma_s^{\Phi_\beta^{\B}}(\pi_\Phi(x))
 =\pi_\Phi(\tau_{-\beta s}(x)).
\tag{4.6}\label{eq:kms-modular}
\]
Rotations preserve the subnet, hence
\[
 \sigma_s^{\Phi_\beta^{\B}}
   (\pi_\Phi(\A(I)))=\pi_\Phi(\A(I)).
\]
Takesaki's theorem now gives the unique
$\Phi_\beta^{\B}$-preserving normal faithful expectation.  By
\eqref{eq:E-preserves}, the group average is such an expectation, so
it is the Takesaki expectation.  Moreover,
\[
 \sigma_s^{\omega_\beta}
 =\sigma_s^{\Phi_\beta^{\B}}|_{\A(I)}
 =\tau_{-\beta s}|_{\A(I)},\qquad
 E_I\sigma_s^{\Phi_\beta^{\B}}
 =\sigma_s^{\omega_\beta}E_I.
\tag{4.7}\label{eq:modular-chain}
\]
Notice the logical order: existence of the Gibbs extension gives a
state on $\B$; modular invariance then identifies its
state-preserving expectation.  Takesaki's criterion alone would not
manufacture a KMS extension from an arbitrary state on $\A$.

\emph{Part (c): removal of the invariance qualifier.}
Let $\psi$ be a rotational $\beta$-KMS state on
$C^*_{\mathrm{ln}}(\B)$ extending $\omega_\beta$.  By
Theorem~\ref{thm:klm} and \eqref{eq:one-sector}, $\B$ has only the
vacuum irreducible DHR sector.  Theorem~\ref{thm:lt} says that every
rotational KMS state is a convex combination of Gibbs states in
irreducible sectors.  There is only one such state,
$\Phi_\beta^{\B}$, so
\[
 \psi=\Phi_\beta^{\B}.
\]
Equation \eqref{eq:gauge-invariance} makes this state
$\Z_4$-invariant, and part (b) then gives
$\psi=\omega_\beta\circ E$.  This proves unconditional uniqueness.
\end{proof}

\section{Where the mixing directions went}

\subsection{Relative-commutant mixing}

Fix an interval and write $N=\A(I)$, $M=\B(I)$.  An internal
multiplicity mixer for the inclusion is an operator $T\in N'\cap M$.
By irreducibility,
\[
 N'\cap M=\C\id,
\tag{5.1}\label{eq:relative-commutant}
\]
so $T=\lambda\id$.  A normalized perturbation cannot use this scalar
to create a new state direction.  Explicitly, if
$T=T^*$ and
\[
 \psi_\varepsilon(x)
 =\frac{\Phi_\beta^{\B}
 ((\id+\varepsilon T)x(\id+\varepsilon T))}
 {\Phi_\beta^{\B}((\id+\varepsilon T)^2)},
\]
then \eqref{eq:relative-commutant} makes numerator and denominator
carry the same scalar factor, and
\[
 \psi_\varepsilon=\Phi_\beta^{\B}.
\tag{5.2}\label{eq:scalar-pert}
\]
Thus every mixing direction that would be implemented by an
operator commuting with the subnet is trivial.

\subsection{Superselection-sector mixing}

A different possible mixing is a convex combination of thermal
states in inequivalent irreducible representations.  KLM gives the
sector count through
\[
 \mu_{\B}=\sum_jd(\rho_j)^2.
\]
Holomorphy makes the sum equal one, so only $\rho_0=\id$ remains.
Longo--Tanimoto identifies the extreme rotational KMS states with
the sector Gibbs states.  Therefore there is no second thermal
sector into which a state can mix.

These are two distinct exclusions:
\begin{center}
\setlength{\fboxsep}{7pt}
\fbox{\begin{minipage}{0.86\textwidth}
\noindent
$\A(I)'\cap\B(I)=\C\id$ removes local multiplicity mixing.

\smallskip\noindent
$\mu_{\B}=1$, together with the Longo--Tanimoto classification,
removes thermal-sector mixing.
\end{minipage}}
\end{center}
Neither exclusion should be substituted for the other.  In
particular, holomorphy alone is a statement about DHR sectors, not
an abstract theorem about all possible dynamics.

\section{TFPT corollary and the P1 reading}

\subsection{The finite matrix shadow}

\begin{corollary}[Interpretation of \texttt{v995}]
\label{cor:v995}
Assume that the TFPT collar scaling limit realizes the net inclusion
of Theorem~\ref{thm:main} with the stated rotation dynamics.  Then
the exact finite results of \texttt{v995} are finite-dimensional
shadows of the two uniqueness mechanisms:
\begin{enumerate}[label=\textup{(\roman*)}]
\item on $\CAR_8=M_{256}$, the KMS equations have rank
  $256^2-1=65535$ on the state tangent and kernel zero;
\item all $33$ selected same-$\Z_4$-charge, unequal-modular-energy
  Hermitian directions fail stationarity and are killed by KMS;
\item on the finite crossed product, fixing the grade-zero
  restriction leaves
  \[
    3\cdot256^2=196608
  \]
  real linear extension coordinates, and invariance under
  $E=(1/4)\sum_r\alpha_r$ kills all of them;
\item covariance is redundant after the full finite-factor KMS
  condition has been imposed.
\end{enumerate}
\end{corollary}

\begin{proof}
Items (i)--(iv) are exactly the finite linear-algebra statements of
\texttt{v995}.  The continuum interpretation is conditional on the
scaling-limit identification because finite matrix rank does not by
itself prove factoriality or classify locally normal KMS states of a
type-III net.  Under that identification, item (iii) is the finite
Fourier version of \eqref{eq:inv-unique}, while items (i), (ii), and
(iv) are the finite CAR instance of Theorem~\ref{thm:araki}.
\end{proof}

\subsection{Equidistribution over four deck channels}

The exact statement needed for the P1 reading is elementary but has
one hypothesis that must be visible.

\begin{lemma}[Transitive four-channel equidistribution]
\label{lem:equidistribution}
Let $p_0,p_1,p_2,p_3$ be orthogonal projections in a response
algebra with $\sum_rp_r=\id$, and suppose a $\Z_4$ action is
transitive:
\[
 \alpha_s(p_r)=p_{r+s\ {\rm mod}\ 4}.
\tag{6.1}\label{eq:transitive}
\]
If $\psi$ is $\Z_4$-invariant, then
\[
 \psi(p_0)=\psi(p_1)=\psi(p_2)=\psi(p_3)=\frac14.
\tag{6.2}\label{eq:quarter}
\]
If a normalized total modular response is
$R_{\mathrm{tot}}=1/\beta_{\mathrm{angle}}$ and is additive over
these channels, then each channel carries
\[
 R_r=\frac1{4\beta_{\mathrm{angle}}}.
\tag{6.3}\label{eq:response-quarter}
\]
\end{lemma}

\begin{proof}
For any $r,s$, invariance and transitivity give
\[
 \psi(p_{r+s})=\psi(\alpha_s(p_r))=\psi(p_r).
\]
All four values are therefore a common number $w$.  Normalization
gives $1=\psi(\id)=\sum_r\psi(p_r)=4w$, hence $w=1/4$.
Additivity of the response then gives
$R_r=R_{\mathrm{tot}}/4$.
\end{proof}

\begin{corollary}[The theorem-conditional P1 identity]
\label{cor:p1}
Under Theorem~\ref{thm:main} and the response-resolution hypothesis
of Lemma~\ref{lem:equidistribution}, the canonical geometric
normalization
\[
 \beta_{\mathrm{angle}}=2\pi
\tag{6.4}\label{eq:beta-angle}
\]
gives
\[
 c_3=R_r
 =\frac1{4\beta_{\mathrm{angle}}}
 =\frac1{4\cdot2\pi}
 =\frac1{8\pi},
\qquad
 4\beta_{\mathrm{angle}}c_3=1.
\tag{6.5}\label{eq:p1}
\]
\end{corollary}

\begin{proof}
Substitute \eqref{eq:beta-angle} into
\eqref{eq:response-quarter}.  No further analytic input is used.
\end{proof}

\begin{remark}[Exactly what P1 gains]
Theorem~\ref{thm:main} proves uniqueness and gauge invariance of the
rotational KMS extension for the abstract holomorphic net.  Once a
four-channel response resolution satisfying
\eqref{eq:transitive} is identified, equality of its weights is no
longer an assumption: it is Theorem~\ref{thm:main} plus
Lemma~\ref{lem:equidistribution}.  The factor $1/4$ is therefore a
theorem-conditional consequence of index four and transitivity.

What is \emph{not} proved is that the TFPT determinant-line response
is represented by projections $p_r$ of this kind, or that its
physical coefficient is the channel response $R_r$.  An outer
fixed-point inclusion of factors does not supply four central local
sector projections: the grading subspaces are not central summands.
That operator identification is the joint
\texttt{P1.INDEX.KMS.01} and
\texttt{ALPHA.QUILLEN.EXACT.01} bridge.
The angular normalization $\beta_{\mathrm{angle}}=2\pi$ is the
geometric Bisognano--Wichmann/modular input, not a consequence of
the group order.  Accordingly, \texttt{AX.P1.01} remains an axiom;
\eqref{eq:p1} is a theorem-conditional identity, not an unconditional
derivation of P1.
\end{remark}

\section{What this closes and does not close}

\begin{description}[style=nextline]
\item[\texttt{SEAM.STATE.RPMIXING.01}, continuum leg.]
For an abstract holomorphic conformal net satisfying the hypotheses,
the rotational KMS extension theorem is complete: there is one
extension, its expectation is canonical, and both local
relative-commutant mixing and thermal-sector mixing are excluded.
The finite \texttt{v995} computation is exactly consistent with this
result.  Applying it to the TFPT collar remains conditional on the
continuum-net identification in the box below.  The ledger's
continuum leg therefore stays open.

\item[\texttt{ALPHA.QUILLEN.EXACT.01}.]
The theorem supplies the rigidity premise: once the U(1)
determinant-line functional is identified with the invariant modular
response of this inclusion, there is no second KMS state or
nontrivial relative-commutant direction that can alter its sector
weights.  It does not prove the Bismut--Freed/Quillen identification,
the zeta-regularized variation, or the compensation normalization.
That contract remains open.

\item[\texttt{AX.P1.01}.]
The arithmetic identity
$1/(4\cdot2\pi)=1/(8\pi)$ and the equidistribution implication are
proved under explicit operator and geometric hypotheses.
The identification of that response with the primitive seam
coefficient is not.  Thus P1 stays typed as an axiom and gains a
theorem-conditional reading, not a marker upgrade.
\end{description}

\section{Remaining gap}

\begin{center}
\setlength{\fboxsep}{10pt}
\fcolorbox{red!65!black}{red!3}{
\begin{minipage}{0.91\textwidth}
\textbf{Remaining gap (exact statement).}
For the TFPT $M=1$ QWZ collar with its four quarter-twisted deck
labels, prove that the scaling-limit local von Neumann algebras are
the hyperfinite type-$\mathrm{III}_1$ factors of a conformal net and
identify, as conformal nets with vacuum, common stress tensor,
rotation dynamics, and $\Z_4$ action, the resulting concrete
inclusion with
\[
 \A=\B^{\Z_4}\subset\B,\qquad
 [\B:\A]=4,\qquad
 \A(I)'\cap\B(I)=\C,\qquad
 \mu_{\B}=1,\quad c_{\B}=8.
\]
The identification must construct the non-bilinear spin/twist
fields, verify relative locality and the commutative index-four
Q-system, and show that the extension is the holomorphic $(E_8)_1$
net rather than merely matching its finite covariance, fusion, and
index arithmetic.  For the P1/Quillen application, the same
identification must additionally show that the normalized
determinant-line modular response is the additive transitive
four-channel response of Lemma~\ref{lem:equidistribution}, with
$\beta_{\mathrm{angle}}=2\pi$.

\medskip
This is the MMST/Osborne--Stottmeister leg beyond fermion bilinears:
the existing theorem controls the bilinear
$\mathfrak{so}(16)_1$ currents, but not the spinor extension fields
that generate $(E_8)_1$.  Finite factoriality
$Z(M_{256})=\C$, the exact covariance \eqref{eq:frozen-C}, the
$33\to0$ KMS rank test, and the $196608\to0$ invariant-extension
test do not prove this type-$\mathrm{III}_1$ scaling-limit
identification.
\end{minipage}}
\end{center}

\section*{Conclusion}
\addcontentsline{toc}{section}{Conclusion}

The abstract operator-algebraic theorem is closed at its natural
strength.  Existence comes from the vacuum-sector rotational Gibbs
state of the holomorphic extension.  The index-four group average is
the unique expectation by the finite-index Radon--Nikodym theorem
and irreducibility.  Takesaki's theorem gives the explicit modular
restriction and intertwining chain.  Gauge-invariant extension
uniqueness is an elementary averaging identity, while unconditional
rotational KMS uniqueness requires both KLM ($\mu=1$ leaves only the
vacuum sector) and Longo--Tanimoto (rotational KMS extreme points are
sector Gibbs states).  Araki's theorem identifies the exact frozen
finite covariance and explains the finite KMS rank collapse.

The remaining issue is not another uniqueness argument.  It is the
construction and identification of the concrete TFPT collar
scaling-limit net, including its non-bilinear extension fields and
its determinant-line response.  That issue is isolated in the
single box above and lies outside Theorem~\ref{thm:main}.

\begin{thebibliography}{99}

\bibitem{Araki1971}
H.~Araki,
\emph{On quasifree states of CAR and Bogoliubov automorphisms},
Publ.\ Res.\ Inst.\ Math.\ Sci.\ \textbf{6} (1970/71), 385--442,
\href{https://doi.org/10.2977/prims/1195193913}
{doi:10.2977/prims/1195193913}.

\bibitem{Boeckenhauer1996}
J.~B\"ockenhauer,
\emph{An algebraic formulation of level one Wess--Zumino--Witten
models},
Rev.\ Math.\ Phys.\ \textbf{8} (1996), 925--948.

\bibitem{BoeckenhauerEvans1998}
J.~B\"ockenhauer and D.~E.~Evans,
\emph{Modular invariants, graphs and $\alpha$-induction for nets of
subfactors I},
Commun.\ Math.\ Phys.\ \textbf{197} (1998), 361--386.

\bibitem{Hiai1991}
F.~Hiai,
\emph{Minimum index for subfactors and entropy II},
J.\ Math.\ Soc.\ Japan \textbf{43} (1991), 347--379.

\bibitem{KLM2001}
Y.~Kawahigashi, R.~Longo and M.~M\"uger,
\emph{Multi-interval subfactors and modularity of representations
in conformal field theory},
Commun.\ Math.\ Phys.\ \textbf{219} (2001), 631--669,
\href{https://doi.org/10.1007/PL00005565}
{doi:10.1007/PL00005565}.

\bibitem{Longo1989}
R.~Longo,
\emph{Index of subfactors and statistics of quantum fields I},
Commun.\ Math.\ Phys.\ \textbf{126} (1989), 217--247.

\bibitem{LongoRehren1995}
R.~Longo and K.-H.~Rehren,
\emph{Nets of subfactors},
Rev.\ Math.\ Phys.\ \textbf{7} (1995), 567--597.

\bibitem{LongoTanimoto2018}
R.~Longo and Y.~Tanimoto,
\emph{Rotational KMS states and type I conformal nets},
Commun.\ Math.\ Phys.\ \textbf{357} (2018), 249--266,
\href{https://doi.org/10.1007/s00220-017-2969-8}
{doi:10.1007/s00220-017-2969-8}.

\bibitem{OsborneStottmeister2023}
T.~J.~Osborne and A.~Stottmeister,
\emph{Conformal field theory from lattice fermions},
Commun.\ Math.\ Phys.\ \textbf{398} (2023), 219--283,
\href{https://doi.org/10.1007/s00220-022-04519-4}
{doi:10.1007/s00220-022-04519-4}.

\bibitem{Takesaki1972}
M.~Takesaki,
\emph{Conditional expectations in von Neumann algebras},
J.\ Funct.\ Anal.\ \textbf{9} (1972), 306--321.

\bibitem{Takesaki2003}
M.~Takesaki,
\emph{Theory of Operator Algebras II},
Encyclopaedia of Mathematical Sciences \textbf{125},
Springer, Berlin, 2003.

\end{thebibliography}

\end{document}
